% UTF-8 LaTeX fragment: reference answers and solutions

\sectionbanner{一、单选题解析}{第1--10题}

\begin{solution}{1}{C}{8.0}{新定义与抽象建模}
\analysis{只需判断奇偶性，将递推转化为奇偶状态的周期。}
\begin{enumerate}[label=\textbf{Step \arabic*:},leftmargin=*]
  \item \(a_1=1\) 为奇数，\(a_2=3\) 为奇数。
  \item 按“后一项等于前两项之和”跟踪奇偶性：
  \[
    \text{奇数}+\text{奇数}=\text{偶数},\qquad
    \text{奇数}+\text{偶数}=\text{奇数},\qquad
    \text{偶数}+\text{奇数}=\text{奇数}.
  \]
  \item 因此奇偶性按
  \[
    \text{奇，奇，偶}
  \]
  循环，每3项中有1项是偶数。
  \item 前30项共有
  \[
    30\div3=10
  \]
  个完整周期，所以偶数项有10项。答案为C。
\end{enumerate}
\methodsummary{递推奇偶题不必求项值，只需记录状态周期。}
\end{solution}

\begin{solution}{2}{C}{6.0}{数论与整除逻辑}
\analysis{三盏灯再次同时亮起的间隔是6、10、15的最小公倍数，再在3小时内计数。}
\begin{enumerate}[label=\textbf{Step \arabic*:},leftmargin=*]
  \item 6、10、15的最小公倍数为
  \[
    \operatorname{lcm}(6,10,15)=30,
  \]
  所以三盏灯每隔30分钟同时亮起一次。
  \item 3小时为
  \[
    3\times60=180
  \]
  分钟。
  \item 不含开始时刻，在180分钟内的同时亮起时刻为
  \[
    30,\ 60,\ 90,\ 120,\ 150,\ 180,
  \]
  共有6次。答案为C。
\end{enumerate}
\methodsummary{周期重合问题用最小公倍数确定共同周期，再注意是否包含起点和终点。}
\end{solution}

\begin{solution}{3}{B}{6.0}{数论与整除逻辑}
\analysis{把最大公因数提出来，转化为两个互素正整数的定和乘积问题。}
\begin{enumerate}[label=\textbf{Step \arabic*:},leftmargin=*]
  \item 因为 \(a,b\) 的最大公因数为12，设
  \[
    a=12x,\qquad b=12y,
  \]
  其中 \(x,y\) 为互素正整数。
  \item 由 \(a+b=60\) 得
  \[
    12x+12y=60,
  \]
  所以
  \[
    x+y=5.
  \]
  \item 互素正整数对可取
  \[
    (x,y)=(1,4),(2,3),(3,2),(4,1),
  \]
  对应 \(xy\) 最大为
  \[
    2\times3=6.
  \]
  \item 因此
  \[
    ab=144xy,
  \]
  最大值为
  \[
    144\times6=864.
  \]
  答案为B。
\end{enumerate}
\methodsummary{最大公因数条件通常先归一化，把变量写成公因数乘互素变量。}
\end{solution}

\begin{solution}{4}{D}{4.0}{几何与空间}
\analysis{\textbf{（分层桥梁法）}破题的核心在于跳出“5个圆依次连成一串”的视觉陷阱，转而将图形按“上层3个圆＋下层2个圆”进行拆解。下层的圆连接着上层相邻的两个圆。仔细观察图2可以看出，\(4\text{ cm}\) 实际上是上下相邻两圆在水平方向上的重叠（或相交）投影跨度。利用下层圆的直径，减去左右两边各自被交叠掉的宽度，即可求出上层两个相邻圆在水平方向上的真正净间距。}
\begin{enumerate}[label=\textbf{Step \arabic*:},leftmargin=*]
  \item 计算上层两圆的净间距。单个下层圆的直径为 \(9\text{ cm}\)。它左侧与上层第1个圆的水平重叠宽度为 \(4\text{ cm}\)，右侧与上层第3个圆的水平重叠宽度同样为 \(4\text{ cm}\)。因此，这个下层圆夹在上层两个圆中间、露出的水平间距（即上层两圆最右缘与最左缘之间的空白距离）为
  \[
    9-4-4=1\text{ cm}.
  \]
  \item 观察图3，决定图案总跨度的是上层排布的3个圆以及它们之间的2处间距。3个圆自身的直径和为
  \[
    3\times9=27\text{ cm},
  \]
  2处相邻圆的间距和为
  \[
    2\times1=2\text{ cm}.
  \]
  所以总长度为
  \[
    27+2=29\text{ cm}.
  \]
\end{enumerate}
\methodsummary{分层降维与投影法（Layered Projection Method）：在处理非单向排列的“上下错落重叠图形”（如五环、蜂窝、砖墙排布）的总长度计算时，切忌用“单列线性平移”的惯性思维去直接乘间距。}
\end{solution}

\begin{solution}{5}{C}{8.0}{排列组合与概率}
\analysis{使用有序结果作为样本空间，先数出乘积能被4整除的条件结果，再数其中点数和为偶数的结果。}
\begin{enumerate}[label=\textbf{Step \arabic*:},leftmargin=*]
  \item 每枚骰子的点数中，奇数有1、3、5，共3个；只含一个因数2的数有2、6，共2个；含两个因数2的数有4，共1个。
  \item 两枚骰子乘积不能被4整除的情形有两类：
  \[
    3\times3=9
  \]
  种两枚都是奇数的结果，以及
  \[
    2\times2\times3=12
  \]
  种一枚为2或6、另一枚为奇数的结果。
  \item 所以乘积能被4整除的有序结果数为
  \[
    36-9-12=15.
  \]
  \item 在这些结果中，点数和为偶数要求两枚骰子同奇偶。由于两枚都是奇数时乘积不能被4整除，所以只需两枚都是偶数。偶数点数为2、4、6，共3个，因此目标结果数为
  \[
    3\times3=9.
  \]
  \item 所求条件概率为
  \[
    \frac9{15}=\frac35.
  \]
  答案为C。
\end{enumerate}
\methodsummary{条件概率题先锁定条件样本空间，再在其中计数目标事件。}
\end{solution}

\begin{solution}{6}{B}{4.0}{几何与空间}
\analysis{利用图形的九折旋转对称性，将三角形按“由外向内、由小到大”进行有序分类，从而将复杂的全局计数转化为“局部数量\(\times9\)＋核心大三角形数”的简明计算。}
\begin{enumerate}[label=\textbf{Step \arabic*:},leftmargin=*]
  \item 外圈小尖角（单块区域）：围绕外圈九个顶点，共9个。
  \item 内圈小尖角（网格交汇）：围绕中心九边形均匀分布，共9个。
  \item 顶角中三角形（跨线拼合）：由外圈顶点延伸到第二条截线形成的组合图，共9个。
  \item 全图大三角形（初始骨架）：贯穿整个圆心区域的最大的正三角形，共3个。
\end{enumerate}
前三种三角形各有9个，一共27个；最后一种有3个。因此总数为
\[
  9+9+9+3=30\text{（个）}.
\]
\methodsummary{面对高阶对称几何图形，应通过观察局部特征“以点带面”，利用对称性成套计数，化繁为简，彻底避免盲目死数带来的遗漏或重复。}
\end{solution}

\begin{solution}{7}{D}{6.0}{新定义与抽象建模}
\analysis{只关心除以7的余数，因此把递推转成余数递推，并寻找周期。}
\begin{enumerate}[label=\textbf{Step \arabic*:},leftmargin=*]
  \item 记 \(r_n\) 为 \(a_n\) 除以7的余数。由递推关系得
  \[
    r_{n+1}\equiv2r_n+1\pmod 7.
  \]
  \item 从 \(r_1=5\) 开始计算余数：
  \[
    r_1=5,\quad r_2=4,\quad r_3=2,\quad r_4=5,
  \]
  所以余数按
  \[
    5,\ 4,\ 2
  \]
  循环。
  \item 因为
  \[
    20\equiv2\pmod3,\qquad21\equiv0\pmod3,
  \]
  所以 \(r_{20}=4,\ r_{21}=2\)。
  \item 因此 \(a_{20}+a_{21}\) 除以7的余数为
  \[
    4+2=6.
  \]
  答案为D。
\end{enumerate}
\methodsummary{递推余数题只跟踪余数状态，找到周期后定位所需项。}
\end{solution}

\begin{solution}{8}{C}{6.0}{几何与空间}
\analysis{时钟夹角问题本质上是“分针追赶时针的环形追及问题”。两次夹角为 \(99^\circ\) 的过程中，分针相比时针多走过的相对路程正好是两个 \(99^\circ\) 之和，即 \(198^\circ\)。用这个角度差除以两针的速度差即可直接得出慢跑时间。}
\begin{enumerate}[label=\textbf{Step \arabic*:},leftmargin=*]
  \item 寻找角速度差。分针每分钟走
  \[
    360^\circ\div60=6^\circ,
  \]
  时针每分钟走
  \[
    30^\circ\div60=0.5^\circ.
  \]
  分针与时针的“速度差”为
  \[
    6^\circ-0.5^\circ=5.5^\circ/\text{分钟}.
  \]
  \item 第一次夹角为 \(99^\circ\) 时，分针在时针“后方” \(99^\circ\)；第二次夹角为 \(99^\circ\) 时，分针在时针“前方” \(99^\circ\)。在这段时间里，分针比时针多走过的总角度差为
  \[
    99^\circ+99^\circ=198^\circ.
  \]
  \item 根据“追及时间＝路程差\(\div\)速度差”，慢跑时间为
  \[
    198^\circ\div\bigl(5.5^\circ/\text{分钟}\bigr)=36\text{分钟}.
  \]
\end{enumerate}
\methodsummary{处理同一小时内两次出现相同夹角的问题，牢记将时钟问题转化为“相对行程追及问题”，利用“两次夹角之和\(\div\)两针速度差＝间隔时间”即可避开繁琐的分步计算。}
\end{solution}

\begin{solution}{9}{C}{8.0}{代数与方程应用}
\analysis{把倒比和通分，转成只含 \(a+b\) 与 \(ab\) 的关系。}
\begin{enumerate}[label=\textbf{Step \arabic*:},leftmargin=*]
  \item 对条件通分：
  \[
    \frac ab+\frac ba=\frac{a^2+b^2}{ab}=\frac52.
  \]
  \item 又因为
  \[
    a^2+b^2=(a+b)^2-2ab,
  \]
  所以
  \[
    \frac{18^2-2ab}{ab}=\frac52.
  \]
  \item 化简得
  \[
    \frac{324}{ab}-2=\frac52,\qquad
    \frac{324}{ab}=\frac92.
  \]
  \item 因此
  \[
    ab=\frac{324\times2}{9}=72.
  \]
  答案为C。
\end{enumerate}
\methodsummary{含 \(\frac ab+\frac ba\) 的题目常转化为 \((a+b)^2\) 与 \(ab\) 的关系。}
\end{solution}

\begin{solution}{10}{C}{8.0}{排列组合与概率}
\analysis{先算全部排列，再利用前两位和与后两位和的对称性，只需扣除相等情形。}
\begin{enumerate}[label=\textbf{Step \arabic*:},leftmargin=*]
  \item 从6个数字中选4个排成一列，总排法为
  \[
    6\times5\times4\times3=360.
  \]
  \item 对于前两个数字之和不等于后两个数字之和的排列，交换前两位组与后两位组，会把“大于”与“小于”一一对应。因此，“大于”的数量是非相等排列的一半。
  \item 计算前两位和等于后两位和的排列。两对数字的和相同的配对为
  \[
  \begin{gathered}
    (1,4),(2,3),\\
    (1,5),(2,4),\\
    (1,6),(2,5),(3,4),\\
    (2,6),(3,5),\\
    (3,6),(4,5).
  \end{gathered}
  \]
  这些可形成7组互不重叠的等和配对。
  \item 每组等和配对中，选择哪一对放前两位有2种，每对内部顺序各有2种，所以每组产生
  \[
    2\times2\times2=8
  \]
  个排列。相等情形共有
  \[
    7\times8=56
  \]
  个。
  \item 所求数量为
  \[
    \frac{360-56}{2}=152.
  \]
  答案为C。
\end{enumerate}
\methodsummary{比较两部分和的排列题可用互补对称，把计数转化为总数减相等情形。}
\end{solution}

\Needspace{8cm}
\sectionbanner{多选题解析}{第11--14题}

\begin{solution}{11}{A、C、D}{6.0}{代数与方程应用}
\analysis{设购买数量，先化简总价方程，再列出满足正整数条件的方案并逐项判断陈述。}
\begin{enumerate}[label=\textbf{Step \arabic*:},leftmargin=*]
  \item 设书签买了 \(s\) 套，彩笔买了 \(c\) 盒。由于两种都买，\(s,c\) 都是正整数。
  \item 由总价得
  \[
    12s+8c=96.
  \]
  两边同除以4，得
  \[
    3s+2c=24.
  \]
  \item 因为 \(2c\) 是偶数，所以 \(3s\) 是偶数，进而 \(s\) 是偶数。
  \item 正整数方案为
  \[
    (s,c)=(2,9),(4,6),(6,3).
  \]
  \item 判断各项：A正确，因为所有方案中 \(s\) 都是偶数；B错误，彩笔盒数只能是9、6、3；C正确，方案 \((4,6)\) 中一共买10件；D正确，书签套数最大为6。所以正确选项为A、C、D。
\end{enumerate}
\methodsummary{购物多选题要先确定全部正整数方案，再分别检验“必然”和“可以”两类陈述。}
\end{solution}

\begin{solution}{12}{A、B、C}{8.0}{数论与整除逻辑}
\analysis{把每个正因数写成 \(2^a3^b5^c\)，再按各选项对指数范围分别计数。}
\begin{enumerate}[label=\textbf{Step \arabic*:},leftmargin=*]
  \item \(N=2^3\times3^2\times5\) 的正因数都可写成
  \[
    2^a3^b5^c,
  \]
  其中
  \[
    0\le a\le3,\qquad0\le b\le2,\qquad0\le c\le1.
  \]
  \item 判断A。正因数个数为
  \[
    (3+1)(2+1)(1+1)=24,
  \]
  A正确。
  \item 判断B。奇正因数要求 \(a=0\)，所以个数为
  \[
    1\times3\times2=6,
  \]
  B正确。
  \item 判断C。能被15整除要求 \(b\ge1\) 且 \(c\ge1\)。因此 \(a\) 有4种选择，\(b\) 有2种选择，\(c\) 有1种选择，个数为
  \[
    4\times2\times1=8,
  \]
  C正确。
  \item 判断D。完全平方数要求各质因数指数均为偶数。此时 \(a\) 可取0或2，\(b\) 可取0或2，\(c\) 只能取0，个数为
  \[
    2\times2\times1=4,
  \]
  不是6，D错误。所以正确选项为A、B、C。
\end{enumerate}
\methodsummary{因数计数题的核心是质因数指数的取值范围，不同限制对应不同指数条件。}
\end{solution}

\begin{solution}{13}{A、B、C}{8.0}{新定义与抽象建模}
\analysis{分别考查总量不变量、奇偶变化和具体操作构造；多选题每个陈述要独立验证。}
\begin{enumerate}[label=\textbf{Step \arabic*:},leftmargin=*]
  \item 每次操作中，一堆增加2枚，另外两堆各减少1枚，总变化为
  \[
    2-1-1=0.
  \]
  所以总数始终不变，A正确。
  \item 一次操作中，被选择的那一堆增加2，奇偶性不变；另外两堆各减少1，奇偶性都会改变。因此每次操作会改变两堆的奇偶性。
  \item 改变两堆奇偶性会使奇数堆的数量变化 \(0,2\) 或 \(-2\)，所以奇数堆数的奇偶性不变。初始三堆为7、10、13，奇数堆有2堆，是偶数。因此之后奇数堆数始终是偶数，B正确。
  \item 给出达到三堆都是10枚的操作：
  \[
    (7,10,13)\longrightarrow(9,9,12)
    \longrightarrow(11,8,11)
    \longrightarrow(10,10,10).
  \]
  前两次选择第一堆操作，第三次选择第二堆操作，所以C正确。
  \item D错误。第一次选择第一堆操作后，第二堆从10变为9，第三堆从13变为12，它们的奇偶性发生改变。所以正确选项为A、B、C。
\end{enumerate}
\methodsummary{操作不变量题要区分总量不变量、奇偶不变量和具体构造，不能只凭一次操作判断全部陈述。}
\end{solution}

\begin{solution}{14}{A、B、C}{8.0}{代数与方程应用}
\analysis{固定和下分别考查乘积极值、奇偶性和整除条件，并用方程关系检验可能值。}
\begin{enumerate}[label=\textbf{Step \arabic*:},leftmargin=*]
  \item 判断A。因为
  \[
    ab\le\left(\frac{a+b}{2}\right)^2=9^2=81,
  \]
  当 \(a=b=9\) 时取到81，所以A正确。
  \item 判断B。由于 \(a+b=18\) 为偶数，\(a,b\) 同奇偶。若二者同为奇数，则 \(a^2,b^2\) 同为奇数，和为偶数；若二者同为偶数，则 \(a^2,b^2\) 同为偶数，和也为偶数。B正确。
  \item 判断C。17是质数。若 \(ab\) 能被17整除，则 \(a\) 或 \(b\) 能被17整除。又 \(a,b\) 是正整数且和为18，所以其中一个数只能是17。C正确。
  \item 判断D。若 \(ab=70\)，则 \(a,b\) 是方程
  \[
    X^2-18X+70=0
  \]
  的两个正整数根。其判别式为
  \[
    18^2-4\times70=44.
  \]
  44不是平方数，所以不存在这样的正整数 \(a,b\)，D错误。所以正确选项为A、B、C。
\end{enumerate}
\methodsummary{固定和多选题要把极值、奇偶和整除分别拆开验证，避免把某个例子误当成全体结论。}
\end{solution}

\sectionbanner{二、填空题解析}{第15--19题}

\begin{solution}{15}{22}{6.0}{排列组合与概率}
\analysis{按除以3的余数分类，和能被3整除只可能来自 \(0+0\) 或 \(1+2\)。}
\begin{enumerate}[label=\textbf{Step \arabic*:},leftmargin=*]
  \item 1到12中除以3余0的数有4个，余1的数有4个，余2的数也有4个。
  \item 两数余数为0和0时，取法为
  \[
    \binom42=\frac{4\times3}{2}=6.
  \]
  \item 两数余数为1和2时，取法为
  \[
    4\times4=16.
  \]
  \item 总取法为
  \[
    6+16=22.
  \]
\end{enumerate}
\methodsummary{余数配对计数要先列出可行余数组合，再分别计算取法。}
\end{solution}

\begin{solution}{16}{106}{6.0}{几何与空间}
\analysis{本题最高效的解法是采用“正难则反”的容斥思维（补集法）：先算出表格中任意选取两个正方形的“总对数”，再减去有公共点的“共边对数（上下左右相邻）”和“共顶点对数（对角相邻）”，剩下的就是没有公共点的正方形对数。}
\begin{enumerate}[label=\textbf{Step \arabic*:},leftmargin=*]
  \item 计算任意选两块的总对数。\(3\times6\) 的表格中一共有
  \[
    3\times6=18
  \]
  个小正方形。从18个正方形中任选2个，总组合对数为
  \[
    \frac{18\times17}{2}=153\text{（对）}.
  \]
  \item 计算“共边（有公共边）”的正方形对数。左右相邻的有
  \[
    (6-1)\times3=15\text{（对）};
  \]
  上下相邻的有
  \[
    (3-1)\times6=12\text{（对）}.
  \]
  共边的总对数为
  \[
    15+12=27\text{（对）}.
  \]
  \item 计算“仅共顶点（对角相邻）”的正方形对数。每个 \(2\times2\) 田字格里恰好有2对对角相邻的正方形，而 \(3\times6\) 表格中包含的 \(2\times2\) 田字格个数为
  \[
    (3-1)(6-1)=2\times5=10.
  \]
  因此仅共顶点的总对数为
  \[
    10\times2=20\text{（对）}.
  \]
  \item 用总数减去共边和共顶点的对数：
  \[
    153-27-20=106\text{（对）}.
  \]
\end{enumerate}
\methodsummary{解决网格图形中的“无重叠/无公共点”计数问题，核心是牢记“正难则反”原则。将“无公共点”转化为“总数－（共边数＋共顶点数）”，能够把复杂的离散空间搜索转化为规则的网格模型计算。}
\end{solution}

\begin{solution}{17}{465}{6.0}{新定义与抽象建模}
\analysis{每次操作把两个数换成它们的和，黑板上所有数的总和保持不变。}
\begin{enumerate}[label=\textbf{Step \arabic*:},leftmargin=*]
  \item 每次擦去 \(a,b\)，再写上 \(a+b\)，黑板上所有数的总和不变。
  \item 因此最后剩下的数等于最初1到30的和。
  \item 1到30的和为
  \[
    \frac{(1+30)\times30}{2}=465.
  \]
  所以最后剩下的数是465。
\end{enumerate}
\methodsummary{操作题先找不变量，本题的不变量是所有数的总和。}
\end{solution}

\begin{solution}{18}{411}{8.0}{几何与空间}
\analysis{本题是小学高年级数学竞赛中经典的“梯形蝴蝶定理”与“一半模型”的综合应用题。破题的关键在于利用正方形左右两对边的平行关系，将两条交叉线看作梯形的对角线，从而构造出“蝴蝶模型”，利用相对三角形面积乘积相等的性质求出隐藏的“上翅膀”面积；再结合正方形的“一半模型”锁定整个正方形的总面积，最后通过简单的容斥减法得出顶部未知区域的面积。}
\begin{enumerate}[label=\textbf{Step \arabic*:},leftmargin=*]
  \item 寻找平行线，构建“梯形蝴蝶模型”。因为正方形的左边缘与右边缘互相平行，夹在左右两边之间的两条交叉长线，刚好可以看作一个平行四边形或梯形内部的两条对角线。对角线相交后，在图中标示的区域内形成了一个标准的“蝴蝶模型”。在这个模型中，左右两个已知面积的三角形（面积分别为350和56）正好对应梯形靠在平行底边上的两块区域，而夹在中间上方、参与构图的三角形就像蝴蝶的“上翅膀”。
  \item 利用“蝴蝶定理”计算“上翅膀”面积。根据梯形蝴蝶定理，对角线分割出的四个三角形中，相对两个三角形的面积乘积相等：
  \[
    \text{左侧面积}\times\text{右侧面积}
    =\text{上翅膀面积}\times\text{下翅膀面积}.
  \]
  由于左右两边缘相互平行，按照几何对称与比例性质，该模型中的“上翅膀”与“下翅膀”面积相等。于是
  \[
    350\times56=19600,\qquad140\times140=19600,
  \]
  所以隐藏的“上翅膀”三角形面积是140。
  \item 利用“一半模型”求正方形总面积。图中由正方形整条左边作为底边、顶点连接到右边缘交点的大三角形，底和高都等于正方形边长，面积是整个正方形面积的一半。该大三角形由左侧已知区域350和刚求出的上翅膀140拼成，因此
  \[
    \text{正方形面积的一半}=350+140=490,
  \]
  正方形的总面积为
  \[
    490\times2=980.
  \]
  \item 用总面积减去题目给出的三块独立区域面积：
  \[
    980-350-163-56=411\text{（平方米）}.
  \]
\end{enumerate}
\methodsummary{解决复杂几何组合图问题的最高效策略是“化零为整，模型叠加”。看到两条长交叉线夹在平行线之间时，优先联想“蝴蝶定理”求过渡面积；看到底边占满正方形单边、顶点触及对边的三角形时，立即联想“一半模型”定全局总和。这种“局部利用蝴蝶定比例\(\to\)局部升维推导总面积\(\to\)整体作差求未知”的三步法，彻底避开了求具体边长与坐标的无解死循环，是解答竞赛压轴题的标准闭环。}
\end{solution}

\begin{solution}{19}{23}{8.0}{代数与方程应用}
\analysis{总和固定，要让最大数尽量小，就让其余数在不超过该最大数的前提下尽量大。}
\begin{enumerate}[label=\textbf{Step \arabic*:},leftmargin=*]
  \item 9个数的总和为
  \[
    9\times18=162.
  \]
  \item 最小数是6，其余8个数都大于6且互不相同。
  \item 若最大数不超过22，则其余8个数的和最大为
  \[
    15+16+17+18+19+20+21+22=148.
  \]
  此时9个数的总和最多为
  \[
    6+148=154,
  \]
  小于162，所以最大数不能不超过22。
  \item 取9个数为
  \[
    6,\ 16,\ 17,\ 18,\ 19,\ 20,\ 21,\ 22,\ 23.
  \]
  它们互不相同，最小数为6，总和为
  \[
    6+16+17+18+19+20+21+22+23=162.
  \]
  \item 因此最大数的最小可能值是23。
\end{enumerate}
\methodsummary{固定总和的最值构造题，要同时证明下界和给出达到下界的一组数。}
\end{solution}

\sectionbanner{三、综合题解析}{第20--22题}

\begin{solution}{20（1）}{21}{4.0}{新定义与抽象建模}
\analysis{灯的状态只有“亮”与“灭”两种，本质是奇偶次叠加（模2加法）。如果一盏灯被触发了奇数次，它最终就亮；被触发了偶数次，状态就相互抵消变回灭。只需要把每次按灯影响的区域看作集合，检查哪些灯被重叠触发两次即可。}
\begin{enumerate}[label=\textbf{Step \arabic*:},leftmargin=*]
  \item 初始所有灯全灭。按1会切换1，2，3，4号灯。
  \item 按5会切换5，7号灯。
  \item 按6会切换6，7号灯。
  \item 7号灯被按5和按6切换了两次，因此最终熄灭；其余1，2，3，4，5，6号灯各被切换一次，最终点亮。数字之和为
  \[
    1+2+3+4+5+6=21.
  \]
\end{enumerate}
\methodsummary{集合异或法（影响范围叠加法）。对于无需考虑顺序的纯操作计算题，不需要一步一步顺着动态过程做，而是直接将每个按动操作构成的“影响子集”全盘合并，进行简单的奇偶次数统计即可破题。}
\end{solution}

\begin{solution}{20（2）}{246}{6.0}{新定义与抽象建模}
\analysis{根据“按下时，只有大于它的相邻数字才变”的规则，一盏灯能够被关掉，能改变它的操作只有它自己，或者比它小的相邻数字。如果当前盘面上最小的亮灯是 \(k\)，那么所有大于 \(k\) 的按键都不可能影响 \(k\)，因此必须直接按 \(k\) 才能把它消掉。}
\begin{enumerate}[label=\textbf{Step \arabic*:},leftmargin=*]
  \item 初始只有2号灯亮着。当前最小的亮灯是2，必须按2。按下2后，2熄灭，4、5亮起。
  \item 当前亮灯为4、5，最小的亮灯是4，必须按4。按下4后，4熄灭，5再次切换而熄灭，6、7亮起。
  \item 当前亮灯为6、7，最小的亮灯是6，必须按6。按下6后，6熄灭，7再次切换而熄灭，所有灯全灭。
\end{enumerate}
总共按了2、4、6，刚好3次。由于顺序无关紧要，任何排列皆可；按题目要求从小到大写为246。
\methodsummary{贪心消元法（自小到大定向清除）。遇到“单向向下传导”的规则时，抓住“最小、定唯一”的破解法则：盘面上现存最小的亮灯是谁，这一步就必须按谁，以此类推产生骨牌效应。}
\end{solution}

\begin{solution}{20（3）}{12357}{6.0}{新定义与抽象建模}
\analysis{完全复用第二问的“定向清除最小数字”方法。}
\begin{enumerate}[label=\textbf{Step \arabic*:},leftmargin=*]
  \item 初始1，4，6号灯亮着。最小亮灯是1，必须按1。1熄灭，2、3亮起，4原本亮着、切换后熄灭，6不受影响。当前亮灯为2，3，6。
  \item 最小亮灯是2，必须按2。2熄灭，4、5亮起。当前亮灯为3，4，5，6。
  \item 最小亮灯是3，必须按3。3熄灭，4原本亮着、切换后熄灭，6原本亮着、切换后熄灭。当前只剩5号灯亮着。
  \item 最小亮灯是5，必须按5。5熄灭，7亮起。
  \item 最后只剩7号灯，必须按7。7熄灭，全部清空。
\end{enumerate}
总共按了1，2，3，5，7号灯，刚好5次。答案为12357。
\methodsummary{状态更新追踪法。这是“单向贪心消元”在多初始点情况下的连续推演。只要每次寻找“当前亮着的灯里数字最小的灯”，按下它并更新整个盘面，循环往复即可得出必然解。}
\end{solution}

\begin{solution}{21（1）}{18}{4.0}{新定义与抽象建模}
\analysis{本题要求找出在1到6的乘法中，恰好能被两对完全不同的因子组成（且因子都不大于6）的乘积。每对不同因子可以通过交换骰子颜色产生2种方法，所以2对因子会产生4种方法。找出这两个乘积6和12并求和即可。}
\begin{enumerate}[label=\textbf{Step \arabic*:},leftmargin=*]
  \item 对两个骰子点数乘积（\(1\times1\) 到 \(6\times6\)，共36种结果）的频数分布进行梳理。设乘积为 \(P\)，统计每个乘积出现的次数。
  \item 出现1次的积共5个：
  \[
    1,\ 9,\ 16,\ 25,\ 36.
  \]
  出现2次的积共10个：
  \[
    2,\ 3,\ 5,\ 8,\ 10,\ 15,\ 18,\ 20,\ 24,\ 30.
  \]
  出现3次的积共1个：
  \[
    4\quad(1\times4,\ 2\times2,\ 4\times1).
  \]
  出现4次的积共2个：
  \[
  \begin{aligned}
    6&:\ 1\times6,\ 2\times3,\ 3\times2,\ 6\times1,\\
    12&:\ 2\times6,\ 3\times4,\ 4\times3,\ 6\times2.
  \end{aligned}
  \]
  验证总数：
  \[
    5\times1+10\times2+1\times3+2\times4=36.
  \]
  \item 出现4次的乘积只有6和12，因此所有符合条件的数值和为
  \[
    6+12=18.
  \]
\end{enumerate}
\methodsummary{乘法表穷举与分类统计法。对于投掷两枚骰子的乘积问题，最稳妥的第一步就是建立“频数分布表”，梳理出哪些乘积出现1次、2次、3次和4次。}
\end{solution}

\begin{solution}{21（2）}{10}{6.0}{新定义与抽象建模}
\analysis{出现2种掷出方法，意味着该乘积只能由唯一的一对不相等的数字相乘得到，例如2对应 \(1\times2\) 和 \(2\times1\)。需要在36种掷法中清点出这类特殊乘积的总个数。}
\begin{enumerate}[label=\textbf{Step \arabic*:},leftmargin=*]
  \item 根据第一问的统计，出现2次的乘积有
  \[
    2,\ 3,\ 5,\ 8,\ 10,\ 15,\ 18,\ 20,\ 24,\ 30.
  \]
  \item 数出这些乘积，共有10个。
\end{enumerate}
\methodsummary{逆向反推法（补集法）。如果第一步已经做好全局统计，可以直接去数。更巧妙的方法是，总共有18个不同的乘积结果，减去出现1次的5个、出现3次的1个、出现4次的2个，剩下的必定是出现2次的10个。}
\end{solution}

\begin{solution}{21（3）}{86}{8.0}{新定义与抽象建模}
\analysis{不要把4个骰子混在一起想，而是分成“红绿”和“蓝黄”两组。如果某一个乘积数值 \(P\)，红绿骰子有 \(n\) 种方法掷出，那么蓝黄骰子同样有 \(n\) 种方法掷出。要让两组的乘积同为 \(P\)，根据乘法原理，方法数是 \(n\times n=n^2\)。将所有可能乘积的 \(n^2\) 加起来就是总答案。}
\begin{enumerate}[label=\textbf{Step \arabic*:},leftmargin=*]
  \item 要求 \(A\times B=C\times D\)，即“红绿骰子的乘积”必须等于“蓝黄骰子的乘积”。
  \item 假设某个乘积数值为 \(P\)，掷出该乘积的方法有 \(n\) 种，那么 \(A\times B=P\) 有 \(n\) 种可能，\(C\times D=P\) 也有 \(n\) 种可能，所以使 \(A\times B=C\times D=P\) 成立的组合有 \(n^2\) 种。
  \item 按频数分类求和：
  \[
  \begin{aligned}
    \text{出现1次的数值（5个）}&:\quad 5\times(1\times1)=5,\\
    \text{出现2次的数值（10个）}&:\quad 10\times(2\times2)=40,\\
    \text{出现3次的数值（1个）}&:\quad 1\times(3\times3)=9,\\
    \text{出现4次的数值（2个）}&:\quad 2\times(4\times4)=32.
  \end{aligned}
  \]
  \item 总组合数为
  \[
    5+40+9+32=86.
  \]
\end{enumerate}
\methodsummary{分组独立与频数平方和定理。处理多变量等式 \(A\times B=C\times D\) 时，将其转化为左右两个独立事件的联合概率模型，把“求总组合数”转化为“求各频数平方之和”。}
\end{solution}

\begin{solution}{22（1）}{153}{6.0}{新定义与抽象建模}
\analysis{本题属于“定点必经路径问题”。因为必须经过展望台（即经过 \(A\) 到 \(B\) 的线段），所以整段路程被截断为两段独立过程：第一段从底面的 \(S\) 爬升到第1层的 \(A\)，第二段从第2层的 \(B\) 爬升到山顶 \(T\)。分别算出两段的走法数量，再用乘法原理相乘。难点在于 \(S\) 和 \(A\) 在同一条垂直棱上，需要考虑“原地上升”和“绕路再上升”两种情况。}
\begin{enumerate}[label=\textbf{Step \arabic*:},leftmargin=*]
  \item 计算 \(S\) 到 \(A\) 的走法。\(S\) 在底面，\(A\) 在 \(\frac13\) 高度的第1层，且 \(A\) 就在 \(S\) 的正上方（同一条棱）。走法可按“在底面找一条棱\(\times\)向上走\(\times\)在第1层绕回 \(A\)”计算。
  \item 不绕路时，底面直接从 \(S\) 上去1种，到第1层就在 \(A\) 点，原地不动1种，共
  \[
    1\times1=1\text{（种）}.
  \]
  \item 绕路时，底面绕到其他4条棱有
  \[
    4\times2=8\text{（种）},
  \]
  上去之后再从该点绕回 \(A\) 有2种，所以共有
  \[
    8\times2=16\text{（种）}.
  \]
  因此
  \[
    [\text{数}]=1+16=17.
  \]
  \item 计算 \(B\) 到 \(T\) 的走法。直接从 \(B\) 所在的棱上去有1种；绕到另外4条棱再上去有
  \[
    4\times2=8\text{（种）}.
  \]
  因此
  \[
    [\text{学}]=1+8=9.
  \]
  \item 由乘法原理，
  \[
    [\text{数}]\times[\text{学}]=17\times9=153.
  \]
\end{enumerate}
\methodsummary{分段乘法与同棱分类法。处理三维空间路径时，先将必经节点作为“截断点”，将大问题拆为小问题。计算单层爬升时，统一采用标准公式：走法＝（底层找棱的方法数）\(\times\)（向上的方法数）\(\times\)（顶层找目标点的方法数）。}
\end{solution}

\begin{solution}{22（2）}{144}{6.0}{新定义与抽象建模}
\analysis{逻辑与第一问一致，分为 \(S\) 到 \(C\) 以及 \(D\) 到 \(T\) 两段。但这一问的难度进阶在于：起点 \(S\) 和终点 \(C\) 不在同一条垂直棱上。必须把5条能够向上的棱分为三类：\(S\) 正上方的棱、\(C\) 正下方的棱，以及其余3条无关的棱。}
\begin{enumerate}[label=\textbf{Step \arabic*:},leftmargin=*]
  \item \(S\) 在底面，\(C\) 在第1层，但 \(C\) 和 \(S\) 不在同一条棱上。按“找棱\(\times\)向上\(\times\)绕到 \(C\)”分类。
  \item 底面直接从 \(S\) 所在的棱上去有1种，第1层绕行到 \(C\) 有2种，共
  \[
    1\times2=2\text{（种）}.
  \]
  \item 底面绕到 \(C\) 正下方的棱有2种，上去之后就在 \(C\) 点、原地不动有1种，共
  \[
    2\times1=2\text{（种）}.
  \]
  \item 底面绕到剩下的3条棱有
  \[
    3\times2=6\text{（种）},
  \]
  上去之后再绕行到 \(C\) 有2种，共
  \[
    6\times2=12\text{（种）}.
  \]
  因此
  \[
    [\text{语}]=2+2+12=16.
  \]
  \item \(D\) 到 \(T\) 的走法与 \(B\) 到 \(T\) 一致，
  \[
    [\text{文}]=9.
  \]
  所以
  \[
    [\text{语}]\times[\text{文}]=16\times9=144.
  \]
\end{enumerate}
\methodsummary{升降电梯分类法。把5条垂直向上的棱看作5部电梯。从 \(S\) 到 \(C\)，只有三种选择：搭乘 \(S\) 门口的电梯（1种走法\(\times\)上面绕2种）、搭乘 \(C\) 门口的电梯（下面绕2种\(\times\)上面不绕1种）、搭乘另外3部远处的电梯（下面绕2种\(\times\)上面绕2种）。这样分类，逻辑清楚严密。}
\end{solution}

\begin{solution}{22（3）}{729}{8.0}{新定义与抽象建模}
\analysis{没有必须经过某点的限制时，实际只面临三个完全独立的跨层动作：底面到第1层、第1层到第2层、第2层到山顶。只需知道任意一次跨层有多少种走法，然后连乘三次即可。}
\begin{enumerate}[label=\textbf{Step \arabic*:},leftmargin=*]
  \item 底面到第1层：寻找向上的棱，有9种走法。
  \item 第1层到第2层：寻找向上的棱，有9种走法。
  \item 第2层到顶点 \(T\)：寻找向上的棱，有9种走法。
  \item 总走法为
  \[
    9\times9\times9=729\text{（种）}.
  \]
\end{enumerate}
\methodsummary{宏观对称与状态转移原理。不管当前站在五边形的哪一个节点，面前永远有5条向上的棱。选择自己脚下这条有1种走法，选择其他4条各有2种走法，即 \(1+4\times2=9\) 种。因为图形完美对称，所以每一次跨越楼层必定产生9种选择。}
\end{solution}
